% Generated draft: requires mathematical and source review.
\documentclass[11pt]{article}
\usepackage[margin=25mm]{geometry}
\usepackage{fontspec}
\setmainfont{Latin Modern Roman}
\usepackage{amsmath,amssymb,mathtools,unicode-math}
\setmathfont{Latin Modern Math}
\usepackage{enumitem,needspace,xurl,hyperref}
\setlength{\parindent}{0pt}
\setlength{\parskip}{5pt}
\setlength{\emergencystretch}{3em}
\newcommand{\N}{\mathbb{N}}
\newcommand{\Z}{\mathbb{Z}}
\newcommand{\Q}{\mathbb{Q}}
\newcommand{\R}{\mathbb{R}}
\newcommand{\C}{\mathbb{C}}
\newcommand{\F}{\mathbb{F}}
\newcommand{\A}{\mathbb{A}}
\newcommand{\PP}{\mathbb{P}}
\newcommand{\E}{\mathbb{E}}
\newcommand{\eps}{\varepsilon}
\newcommand{\dd}{\,\mathrm{d}}
\begin{document}
\section*{UnsolvedMath ID: 20002377}
\textbf{Source problem code:} catalogue entry
\subsection{Definitions and mathematical statement}
\textbf{Source status.} Item A.5 is Natasha Dobrinen's participant-interest statement, not a problem with prescribed hypotheses and a requested conclusion. It expresses a wish to learn core-model and large-cardinal techniques for studying relationships among Boolean-algebra games, distributive laws and stationary subsets of \(P_\kappa\lambda\). No particular game, distributive law, large-cardinal assumption or equivalence is posed. Accordingly, the following records the aid's mathematical synthesis separately, rather than inventing a theorem as the workshop question.

For a regular uncountable cardinal \(\kappa\) and an ordinal \(\lambda\geq\kappa\), put
\[
 P_\kappa(\lambda)=\{x\subseteq\lambda:|x|<\kappa\}.
\]
A club in this set is unbounded under inclusion and closed under increasing unions of nonempty chains of length less than \(\kappa\); stationary sets meet every club. For forcing, stronger conditions are written smaller. A partial order is \(<\kappa\)-closed if every decreasing sequence of length less than \(\kappa\) has a lower bound. Use \(<\kappa\)-distributive to mean that the intersection of any family of fewer than \(\kappa\) dense open subsets is dense. For a complete Boolean algebra this refers to its nonzero part as a forcing order. Countable distributivity here includes countable families, hence corresponds to \(<\omega_1\)-distributivity.

\textbf{Reported game implication, with missing definition.} The research aid states that \(<\kappa\)-closure implies a winning strategy in an explicitly defined dense-set fusion game, and that this winning property implies \(<\kappa\)-distributivity. However, the supplied summary contains none of that game's rules: its length, move order, legality at limits and winning condition are absent. Denoting that unspecified winning property by \(\mathcal W_\kappa\), the reported schematic implication is
\[
 <\kappa\text{-closure}\ \Longrightarrow\ \mathcal W_\kappa
 \ \Longrightarrow\ <\kappa\text{-distributivity}.
\]
This is not a fully formulated game theorem until the omitted definition is recovered; no alternative standard game is silently substituted, and strictness of either arrow is not established by the supplied details.

\textbf{Reported failure of a stationarity-preservation extension.} For every cardinal \(\lambda\geq\omega_1\), the aid reports a countably distributive forcing Boolean algebra destroying a ground-model stationary subset of \(P_{\omega_1}(\lambda)\). Its described club-shooting construction can be stated with the necessary stationary-complement qualification as follows. In the ground model choose a stationary costationary \(A\subseteq\omega_1\), let \(S=\omega_1\setminus A\), and let \(\mathbb P_S\) consist of countable closed bounded subsets of \(S\), ordered by end extension. Let \(\mathbb B_S=\operatorname{RO}(\mathbb P_S)\) be its regular-open completion. Put
\[
 T_{A,\lambda}=\{x\in P_{\omega_1}(\lambda):
 x\cap\omega_1\text{ is an ordinal belonging to }A\}.
\]
The aid reports that \(T_{A,\lambda}\) is stationary in the ground model, that \(\mathbb B_S\) is countably distributive, and that forcing with it adds a club \(C\subseteq S\), making \(T_{A,\lambda}\) nonstationary. The intended disjoint club is obtained by requiring \(x\cap\omega_1\) to be an ordinal in \(C\). For this standard club-shooting argument, stationarity of \(S\) is essential to the distributivity claim; the summary's bare assumption that \(A\) is stationary is insufficient without this qualification. Thus countable distributivity alone does not ensure preservation of all such stationary sets. This counterexample does not assert that every stationary set is destroyed, nor that closure or a fully specified stronger game property fails to preserve stationarity.

The source's research program remains context rather than an unresolved yes/no assertion, and the aid's later-paper and counterexample claims require independent source verification.
\subsection{Short English statement}
Dobrinen's item describes a research program about games, distributivity and stationarity, rather than posing a specific problem. The aid adds an incompletely specified game implication and a club-shooting example showing that countable distributivity does not by itself preserve stationary subsets of \(P_{\omega_1}(\lambda)\). The necessary stationary-complement assumption is made explicit.
\subsection{Sources}
\begin{itemize}
\item Local UnsolvedMath registry, permanent ID \(20002377\), catalogue entry, \texttt{aim-logic-notes.json}, \texttt{notes[152]}, workshop \emph{Recent advances in core model theory}, participant section A.5, Natasha Dobrinen: \url{https://aimath.org/WWN/coremodel/coremodel.pdf}.
\item Supplied research summary on a dense-set fusion game and club shooting, and its assertion that Dobrinen's later primary papers substantially realized the program. No paper titles, bibliographic references or full game definition are supplied; no independent literature assessment is made here.
\item Attempt 1 is labelled context\_only, problem status at run not\_a\_problem, rigor ``proved'', low novelty confidence. The registry separately says open/OPEN-TRIAGE and statement status unrecoverable. These metadata do not create a mathematical question. The machine-generated research aid is not peer reviewed and requires independent verification before citation.
\end{itemize}
\end{document}
